% Propietario conservado: /Users/ruben/Documents/New project/output/REVISION_ARGUMENTAL_APERTURAS_CIERRES_20260915/VI/delivery/MANUSCRIPT_VI_ES_EN_COMPLETE/source_es/sections/14_generacion_regional.tex
\section{Conditions initiales et génération régionale}
\label{vi:vi:reg:seccion}\label{vi:sec:extension}\label{vi:sec:generacion}

Une position d’APP ne contient pas à elle seule la condition initiale de sa
dynamique. Celle-ci inclut la direction de transport et distingue les curseurs
des deux feuillets. Avant de classer les particules, il est nécessaire de
construire ce domaine : c’est de lui que proviennent les registres régionaux qui
interviendront dans la coordonnée électronique, l’action et les lecteurs de
masse. Les opérations de cette section ne reçoivent ni espèce physique ni
valeur métrologique comme condition initiale.

\subsection{Deux curseurs orientés sur le support complet}

Nous employons les indices \(I_9=\{0,\ldots,8\}\), dont l’étiquette positive est
\(i+1\), et les directions cardinales \(\mathcal D=\{N,E,S,O\}\).
Le domaine d’un curseur et le domaine conjoint sont
\begin{equation}
 \mathcal S=I_9^2\times\mathcal D,\qquad
 \Omega_{\mathrm{seed}}=\mathcal S_+\times\mathcal S_\times,
 \qquad |\mathcal S|=324,\quad
 |\Omega_{\mathrm{seed}}|=104\,976.
 \label{vi:eq:semillas}
\end{equation}
L’indice \(\mathrm{seed}\) conserve la dénomination du domaine dans
les antécédents ; mathématiquement, il s’agit de conditions initiales
orientées. Le produit cartésien inclut toutes les paires ordonnées.
Le feuillet additif dispose d’un curseur et le multiplicatif d’un autre ; une
coïncidence ultérieure entre leurs lectures est un résultat du parcours.

Pendant les \(54\) instants de l’émission initiale, le sélecteur de phase
répète \((+1,-1,0)\). Dans le régime positif, on évalue d’abord
\(\rho_9(i+j+2)\) au curseur additif, puis on transporte ce
curseur d’une position. Dans le négatif, on évalue
\(\rho_9((i+1)(j+1))\) au multiplicatif, puis on transporte
le second curseur. Le régime nul enregistre un événement neutre. Les
directions sont inversées après les instants \(27\) et \(54\).
Cette dernière inversion appartient à l’état terminal bien qu’elle ne modifie pas
les lectures déjà effectuées.

L'observable multiplicatif utilise la carte des exposants du groupe des
unités de \(\ZZ/9\ZZ\) :
\[
 \ell_9(1)=0,\quad\ell_9(2)=1,\quad\ell_9(4)=2,\quad
 \ell_9(8)=3,\quad\ell_9(7)=4,\quad\ell_9(5)=5.
\]
Dans le secteur \(\{3,6,9\}\), cette observable prend la valeur zéro.
L’appartenance au secteur radical reste présente dans l’état ; attribuer
cette valeur à une lecture scalaire n’identifie pas le radical à l’unité
multiplicative.

\subsection{Formule finie des signatures de fenêtre}

Les six fenêtres ont neuf instants chacune. Chaque feuillet est évalué
trois fois par fenêtre. Une formule fermée permet de reconstruire les
signatures sans consulter de liste préalable. Pour un curseur d’origine
\(p=(i,j)\) et de direction \(d\), soit
\[
 f(k)=\begin{cases}k,&0\le k<9,\\18-k,&9\le k<18,
 \end{cases}\qquad p_k=p+f(k)\nu_d\pmod9.
\]
L’indice \(k\) compte seulement les évaluations de ce feuillet.
L’inversion après neuf de ses évaluations explique la seconde
expression de \(f\). En écrivant \(p_k=(i_k,j_k)\), pour
\(m=0,\ldots,5\), on obtient
\begin{align}
 S_m(p,d)&=\sum_{k=3m}^{3m+2}\rho_9(i_k+j_k+2),\notag\\
 E_m(p,d)&=\sum_{k=3m}^{3m+2}
       \ell_9\!\left(\rho_9((i_k+1)(j_k+1))\right),\notag\\
 s_m&=[S_m]_{10},\qquad e_m=[E_m]_{10}.
 \label{vi:vi:reg:firmas}
\end{align}
L’évaluation entière et son quotient de division restent stockés dans
le registre. Les signatures précédentes sont les projections décimales
spécifiées par l’émetteur.

Comme chaque fenêtre comporte \(Z_m=3\) événements neutres, la règle qui
couple les deux signatures est
\begin{align}
 d_{m,1}&=[S_m]_{10},\qquad
 d_{m,2}=[S_m+E_m]_{10},\notag\\
 d_{m,3}&=[3S_m+5E_m+7Z_m]_{10},\notag\\
 U_m&=100s_m+10[s_m+e_m]_{10}+[3s_m+5e_m+1]_{10}.
 \label{vi:eq:emisor-seis}
\end{align}
Le terme final \(1\) est la réduction de \(7Z_m=21\) modulo
dix. Les coefficients et le calendrier sont déclarés dans la loi
d’émission ; les grandeurs physiques n’apparaissent pas parmi ses arguments.
La projection ternaire orientée est
\begin{equation}
 \Psi(U_0,\ldots,U_5)=([-U_0]_3,\ldots,[-U_5]_3)
 \in\FF_3^6.
 \label{vi:eq:psi-catalogo}
\end{equation}
L'identification équilibrée du TRIT est ensuite appliquée. Le signe
moins appartient à cette carte et doit être transporté lors de son changement.

\begin{proposition}[Factorisation et recensement de l’émission]
\label{vi:vi:reg:censo}
Les conditions initiales produisent \(18\) signatures additives et \(26\)
multiplicatives. Leur couplage est injectif et comporte \(468\)
émissions. L’image de \(\Psi\) contient \(243\) mots dans
l’espace ambiant à \(729\) éléments.
\end{proposition}
\begin{proof}
Les centaines de \(U_m\) récupèrent \(s_m\). Si \(b_m\) est le
chiffre des dizaines, \(e_m=[b_m-s_m]_{10}\) ; le chiffre des unités
vérifie la troisième congruence. Ainsi, aucune paire distincte
de signatures n’a la même émission.

Pour obtenir le recensement, on évalue \eqref{vi:vi:reg:firmas} en
\(i,j=0,\ldots,8\) et dans les quatre directions. L’annexe
\ref{vi:vi:reg:tablas} contient toutes les signatures résultantes. Chaque
signature additive possède \(18\) antécédents. Parmi les multiplicatives,
\(24\) en possèdent \(8\), une en possède \(24\) et une en possède \(108\).
Ces sommes sont respectivement \(18\cdot18=324\) et
\(24\cdot8+24+108=324\). Puisque les deux curseurs sont
indépendants, les multiplicités se multiplient. Il en résulte
\(432\) émissions de poids \(144\), \(18\) de poids \(432\) et
\(18\) de poids \(1944\), de somme
\[
 432\cdot144+18\cdot432+18\cdot1944=104\,976.
\]
Les six réductions de \eqref{vi:eq:psi-catalogo} sont appliquées à chaque paire
des \(18\) et \(26\) signatures publiées dans l'annexe. La
décomposition complète en orbites qui y figure donne
\(38\cdot6+5\cdot3=243\) images différentes. Toutes les
opérations de cette énumération sont les sommes, produits et résidus
finis des formules précédentes.
\end{proof}

La démonstration finie n’identifie pas l’exhaustivité du domaine à une
profondeur illimitée. Les conditions initiales ont été recensées sur
un support fixe. La profondeur est construite par des prolongements
compatibles qui conservent ces conditions et leurs transports.

\subsection{Sélection des régions par fibres orientées}

Pour \(w=(w_1,\ldots,w_6)\), définissons
\[
 r(w)=(w_3,w_1,w_2,w_5,w_6,w_4),\qquad
 s(w)=(w_4,w_5,w_6,w_1,w_2,w_3).
\]
Les identités \(r^3=s^2=I\), \(srs=r^{-1}\) réalisent
\(D_3\). Permuter de la même manière les six composantes de \(U\)
commute avec \(\Psi\). L’énumération précédente vérifie que les
orbites conservent les multiplicités de leurs fibres.

Soient \(f(w)\) le nombre d’émissions sur \(w\), \(m(w)\)
leur multiplicité totale et \(\nu(w)=(n_0,n_1,n_2)\) leur recensement
ternaire. La signature additive conserve la classe diagonale
\(d=[i+j]_9\) de l’origine et le sens \(ES\) ou \(NO\).
Nous écrivons \(\mathcal D(\mathcal O)\) pour l’ensemble de ces
classes sur une orbite.

\begin{proposition}[Régions distinguées du catalogue]
\label{vi:vi:reg:seleccion}
Il existe une unique orbite de six mots avec cinq émissions par
mot et des multiplicités \(432+4\cdot144\). Parmi les orbites
de six mots qui satisfont
\[
 f(w)=1,\quad m(w)=144\quad(w\in\mathcal O),\qquad
 \mathcal D(\mathcal O)=\{0,3,6\},
\]
le recensement \((2,3,1)\) et le recensement \((2,2,2)\) sélectionnent,
chacun, une unique orbite. La marque \(d=6\) avec l’orientation
\(ES\) détermine les représentants
\[
 w^{\rm cl}=010211,\qquad w^{\rm pr}=201101,
 \qquad w^{\rm au}=121200.
\]
\end{proposition}
\begin{proof}
Le tableau complet des \(43\) orbites dans
\ref{vi:vi:reg:tablas} permet d’appliquer les prédicats. L’orbite de
clôture est la seule ayant une fibre de taille cinq et de poids \(1008\) ;
ses cinq poids sont ceux indiqués. La condition diagonale et la
fibre simple de poids \(144\) laissent six orbites. Leurs recensements
sélectionnent respectivement une orbite avec \((2,3,1)\) et une avec
\((2,2,2)\). Appliquer les six permutations de \(D_3\) et
vérifier la classe diagonale et l’orientation dans
\eqref{vi:vi:reg:firmas} fixe les trois représentants écrits.
L’unicité utilise toutes les conditions : si l’on omet l’ensemble des
diagonales, quatre orbites équilibrées de fibre simple
et de poids \(144\) apparaissent.
\end{proof}

Les désignations clôture, propagation et auto-échelle s’appliquent
d’abord à ces régions. La section suivante introduit leurs
caractères d’évaluation. En particulier, la signature multiplicative
\(555555\) distingue une émission transversale de clôture, tandis
que \((5,5)\) est une position du support électronique. Les deux
constructions sont reliées par des opérations ultérieures et non
par l’identité de leurs chiffres.

\subsubsection{Involution spéculaire des composantes régionales}
\label{vi:mono-ampl:regiones}

Sur le bloc régional
\(B=(u^{\rm cl},u^{\rm pr},u^{\rm au})\in(\FF_3^6)^3\),
l’échange est
\[
 \mathfrak cB=(u^{\rm pr},u^{\rm cl},u^{\rm au}).
\]
Son carré est l’identité et il maintient la composante d’auto-échelle.
Si \(q=u^{\rm cl}+u^{\rm pr}+u^{\rm au}\),
\(a=u^{\rm cl}+u^{\rm pr}-u^{\rm au}\) et
\(d=u^{\rm cl}-u^{\rm pr}\), on récupère le bloc par
\[
 u^{\rm au}=2(q-a),\quad s=q-u^{\rm au},\quad
 u^{\rm cl}=2(s+d),\quad u^{\rm pr}=2(s-d).
\]
La démonstration utilise \(2^{-1}=2\) dans \(\FF_3\). L’échange
conserve \(q,a,s\) et change le signe de \(d\). Par conséquent,
l’agrégat symétrique et la différence orientée conservent ensemble
les deux composantes que l’agrégat isolé ne distingue pas.
Cette involution régionale et l’inversion d’une route cardinale
agissent sur des domaines distincts. La composante fixe peut
évoluer pendant le transport tout en étant fixe sous l’échange.

\subsection{Calendrier, monodromie et mémoire}

Le calendrier fini se compose de douze fenêtres nonadiques. Sa
suite de groupes est
\begin{equation}
 0\longrightarrow C_{12}\xrightarrow{[b]\mapsto[9b]}C_{108}
 \xrightarrow{[t]\mapsto[t]_9}C_9\longrightarrow0.
 \label{vi:eq:extension108}
\end{equation}
Pour \(t=9b+r\), \(0\le r<9\), la composition contient
le quotient de transport
\[
 (b,r)+(b',r')=
 \left(b+b'+\left\lfloor\frac{r+r'}9\right\rfloor,
       [r+r']_9\right),
\]
en réduisant la première coordonnée modulo douze. La projection est
exacte car son noyau est constitué des multiples de neuf. Elle n’admet pas de
section homomorphe : celle-ci donnerait \(C_{108}\cong C_{12}\times C_9\),
mais les exposants respectifs sont \(108\) et \(36\).

La monodromie décrit l’action d’un tour de la base de phase
sur sa fibre. Dans la dynamique avec mémoire, l’holonomie nonadique
\(\Gamma_9\) conserve la phase et avance la profondeur. Sa
projection sur la phase et le compteur vérifie
\[
 (r,n)\longmapsto(r,n+1).
\]
Cette formule est une représentation réduite de l’action, non une
définition qui élimine le reste de l’état. Sur des histoires
bidirectionnelles, le compteur peut être étendu à \(\ZZ\) ; sur des
histoires commencées à profondeur zéro, le recul n’est
défini que là où le prédécesseur est conservé. Le calendrier peut
revenir au bout de \(108\) transitions et conserver une mémoire
distincte de l’initiale. Cette distinction sera essentielle lorsque la
particule sera construite comme persistance et non comme répétition d’une
coordonnée visible.

\subsection{Des registres aux opérateurs de prolongement}

Les opérateurs linéaires de publication sont déterminés sur les
transitions du transport. Si \(x_j^\chi\) est un état régional
et \(\Pi_6\) sa publication à six composantes, la dépendance est
\[
 b_j^\chi=\Pi_6(x_j^\chi),\qquad
 x_{j+1}^\chi=U_{9j+9}\cdots U_{9j+1}(x_j^\chi).
\]
Les matrices des origines \(X_a\) et des destinations \(Y_a\), en
vecteurs lignes de \(\FF_3^6\), sont
\[
\begin{array}{c|c|c|c}
X_0&Y_0&X_1&Y_1\\\hline
010211&012222&010211&002111\\
012222&010211&002111&110221\\
201101&121221&102011&012222\\
121221&102011&012222&102011\\
121200&112202&121020&010210\\
112202&121020&010210&010200
\end{array}
\]
On a \(\det X_0=1\), \(\det X_1=-1\) dans \(\FF_3\).
La solution unique de \(X_aL_a=Y_a\) est \(L_a=X_a^{-1}Y_a\).
La concaténation des blocs publiés dans le calendrier
\((L_0,L_0,L_1,L_1)\) produit des longueurs
\(6,12,18,24,30\). La frontière de longueur \(36\) conserve
le terminal et la relation de prolongement. Ces matrices
représentent des transitions : l'inversion linéaire prouve leur unicité
pour le registre donné, tandis que la génération des états
appartient à la composition effective \(U_t\) du TPK.

\subsection{Portée de la construction régionale dans l’étude des masses}

Le recensement a séparé conditions initiales, émissions, régions et
opérateurs représentant leurs transitions. L’étape suivante
évalue les coordonnées régionales et compose le registre dodécaphasique.
Plus loin, l’atlas des particules sera construit sur le même
support mais utilisera une autre classification. Cette continuité préserve
une provenance commune sans transformer une fibre numérique en une
espèce par simple coïncidence de cardinaux ou d’étiquettes.
